\section{表与图}这节用来展示表格与图片的插入。

\subsection{表格}
\par 本来LaTeX里表格的变化是非常多的，但鉴于学校要求用三线式，如表\ref{table:example}所示，问题反而简单了。以下是一个例子：
\begin{table}[htbp]\center\songti\zihao{5}
    \caption{示例表格\\Example Table}
         \label{table:example}
    \begin{tabular}{lcccccl}
     \toprule
     例子 & 。。 & 123。 & 。。 & 。。& 。。 & 。。\\
     \midrule
    。。 & 。。 & 。。 & 。。 & 。。& 。。 & 。。\\
    。。 & 。。 & 测试 & 。。 & 。。& 。。 & 。。\\
    。。 & 。。 & 。。 & 。。 & 。。& 。。 & 。。\\
    。。 & 。。 & 。。 & 。。 & 。。& 。。 & 。。\\
    。。 & 。。 & 。。 & 。。 & 。。& 。。 & 。。\\
     \bottomrule
    \end{tabular}
   \end{table}
如果你有使用更复杂的表格的需求，请自行查资料完成。

\subsection{插图}
由于这份模板不考虑多栏排版，所以格式要求中所属的半栏图大小要求我们不作演示。以下是一个通栏图的演示：(如图\ref{fig:logo}所示)
%\begin{figure}[H]
%    \centering
%    \includegraphics[width=100mm]{./figures/figure_example1.jpg}
%    \caption{图片测试（最小宽度）\\Image test (Minimal width)}
%  \end{figure}

\begin{figure}[H]
	\centering
	\includegraphics[width=0.3\linewidth]{./figures/figure_example1.jpg}
	\caption{ElegantLaTeX Logo}
	\label{fig:logo}
\end{figure}


%\begin{figure}[H]
%    \centering
%    \includegraphics[width=130mm]{./figures/figure_example1.jpg}
%    \caption{图片测试（最大宽度）\\Image test (Maximal width)}
%\end{figure}
\par 注意：这里为了减少图片上下的空白，使用了float宏包。